\documentclass{article}
\usepackage{amsmath,amssymb,amsthm}
\usepackage{algorithm}
\usepackage{algpseudocode}
\usepackage{bm}
\usepackage{hyperref}
\usepackage{enumitem}
\newtheorem{theorem}{Theorem}
\newtheorem{lemma}{Lemma}
\newtheorem{assumption}{Assumption}
\newcommand{\E}{\mathbb{E}}
\newcommand{\R}{\mathbb{R}}
\newcommand{\norm}[1]{\left\lVert #1 \right\rVert}
\newcommand{\inner}[2]{\left\langle #1, #2 \right\rangle}
\newcommand{\grad}{\nabla}
\begin{document}

\section*{Algorithm Name}
\textbf{Local Stochastic Gradient Descent (Local SGD) / Federated Averaging (FedAvg)}  

The method performs $E$ local SGD updates on each of $K$ clients before a global averaging step.

\section*{Mathematical Formulation}
Let $f_i:\R^d\to\R$ be the local objective on client~$i$,
\[
f_i(w)=\E_{z\sim\mathcal D_i}\bigl[ \ell(w;z)\bigr],\qquad 
f(w)=\frac1K\sum_{i=1}^K f_i(w).
\]
Denote by $w_t\in\R^d$ the \emph{global} model after the $t$‑th communication round and by
$w_{t}^{i,\,e}$ the local copy on client $i$ after $e\in\{0,\dots,E\}$ local steps, where
$w_{t}^{i,0}=w_t$.

For a chosen stepsize $\eta_t>0$, each client performs $E$ stochastic gradient steps:
\[
\boxed{%
w_{t}^{i,\,e+1}=w_{t}^{i,\,e}-\eta_t g_{t}^{i,\,e},\qquad 
g_{t}^{i,\,e}= \grad\ell\bigl(w_{t}^{i,\,e};\zeta_{t}^{i,\,e}\bigr)
}
\tag{1}
\]
where $\zeta_{t}^{i,\,e}\sim\mathcal D_i$ is an i.i.d. sample on client~$i$.
After $E$ local steps the server aggregates:
\[
\boxed{%
w_{t+1}= \frac1K\sum_{i=1}^K w_{t}^{i,E}
}
\tag{2}
\]

Define the \emph{client drift} (the quantity that appears because local models move away from the global model):
\[
\delta_t \;:=\; \frac1K\sum_{i=1}^K \bigl(\grad f_i(w_t)-\grad f(w_t)\bigr).
\tag{3}
\]

\section*{Pseudocode}
\begin{algorithm}[H]
\caption{Local SGD (FedAvg)}\label{alg:localsgd}
\begin{algorithmic}[1]
\State \textbf{Input:} stepsizes $\{\eta_t\}_{t\ge0}$, local steps $E$, number of clients $K$,
initial model $w_0\in\R^d$.
\For{$t=0,1,\dots,T-1$}
    \State Server broadcasts $w_t$ to all clients.
    \For{each client $i\in\{1,\dots,K\}$ \textbf{in parallel}}
        \State $w_{t}^{i,0}\gets w_t$
        \For{$e=0$ \textbf{to} $E-1$}
            \State Sample $\zeta_{t}^{i,e}\sim\mathcal D_i$
            \State $g_{t}^{i,e}\gets \grad\ell\!\bigl(w_{t}^{i,e};\zeta_{t}^{i,e}\bigr)$
            \State $w_{t}^{i,e+1}\gets w_{t}^{i,e}-\eta_t g_{t}^{i,e}$ \Comment{local SGD step}
        \EndFor
    \EndFor
    \State Server aggregates: $w_{t+1}\gets \frac1K\sum_{i=1}^K w_{t}^{i,E}$
\EndFor
\State \textbf{Output:} $w_T$
\end{algorithmic}
\end{algorithm}

\section*{Dry Run Example}
Consider a single‑dimensional quadratic $f(w)=(w-3)^2$.
Take $K=2$ clients, $E=1$, stepsize $\eta=0.1$, and start from $w_0=0$.
Assume both clients see the same deterministic gradient (no stochasticity).

\begin{enumerate}[label=\arabic*.]
\item \textbf{Round $t=0$:}
\begin{align*}
\grad f_i(w_0)&=2(w_0-3) = -6,\\
w_{0}^{i,1}&=w_{0}^{i,0}-\eta\,\grad f_i(w_{0}^{i,0})
            =0-0.1(-6)=0.6 .
\end{align*}
Aggregation:
\[
w_1=\frac{1}{2}(0.6+0.6)=0.6 .
\]
Objective value: $f(w_1)=(0.6-3)^2=5.76$.

\item \textbf{Round $t=1$:}
\[
\grad f_i(w_1)=2(0.6-3)=-4.8,\qquad
w_{1}^{i,1}=0.6-0.1(-4.8)=1.08 .
\]
Aggregate:
\[
w_2=1.08,\qquad f(w_2)=(1.08-3)^2=3.6864 .
\]

\item \textbf{Round $t=2$:}
\[
\grad f_i(w_2)=2(1.08-3)=-3.84,\qquad
w_{2}^{i,1}=1.08-0.1(-3.84)=1.464 .
\]
Aggregate:
\[
w_3=1.464,\qquad f(w_3)=(1.464-3)^2=2.3645 .
\]
\end{enumerate}
The iterates move toward the optimum $w^\star=3$ while the objective decreases at each round.

\newpage
\section*{Assumptions}
\begin{assumption}[Convexity and Smoothness]\label{as:convexsmooth}
Each local loss $f_i$ is convex and $L$‑smooth:
\[
\forall w,w'\in\R^d,\qquad
f_i(w')\le f_i(w)+\inner{\grad f_i(w)}{w'-w}+\frac{L}{2}\norm{w'-w}^2 .
\]
Consequently $f$ is also convex and $L$‑smooth.
\end{assumption}

\begin{assumption}[Bounded Stochastic Gradient Variance]\label{as:variance}
For any $i$ and any $w$,
\[
\E_{\zeta\sim\mathcal D_i}\!\bigl[\norm{ \grad\ell(w;\zeta)-\grad f_i(w)}^2\bigr]\le\sigma^2 .
\]
\end{assumption}

\begin{assumption}[Bounded Gradient Norm]\label{as:boundedgrad}
\[
\forall i,w:\qquad \norm{\grad f_i(w)}\le G .
\]
\end{assumption}

\begin{assumption}[Bounded Client Drift]\label{as:drift}
Define the drift $\delta_t$ in \eqref{3}. There exists a constant $D\ge0$ such that
\[
\norm{\delta_t}\le D,\qquad\forall t\ge0 .
\]
\end{assumption}

\begin{assumption}[Stepsize]\label{as:stepsize}
The stepsize is non‑increasing and satisfies
\[
\eta_t \le \frac{1}{2L},\qquad
\sum_{t=0}^{\infty}\eta_t =\infty,\qquad
\sum_{t=0}^{\infty}\eta_t^2 <\infty .
\]
A typical choice is $\eta_t=\frac{\eta_0}{\sqrt{t+1}}$.
\end{assumption}

\section*{Main Convergence Theorem}
\begin{theorem}[Convergence of Local SGD]\label{thm:main}
Under Assumptions \ref{as:convexsmooth}–\ref{as:stepsize},
let $w_T$ be the model after $T$ communication rounds.
Define the averaged iterate $\bar w_T:=\frac{1}{T}\sum_{t=0}^{T-1} w_t$.
Then
\[
\E\!\bigl[f(\bar w_T)-f(w^\star)\bigr]
\;\le\;
\frac{\norm{w_0-w^\star}^2}{2\eta_0\sqrt{T}}
\;+\;
\frac{L\eta_0\sigma^2}{K\sqrt{T}}
\;+\;
\frac{L E^2 \eta_0 D^2}{\sqrt{T}} .
\tag{4}
\]
Consequently $\E[f(\bar w_T)-f(w^\star)]=\mathcal O\!\bigl(T^{-1/2}\bigr)$.
\end{theorem}

\section*{Theoretical Properties}
\begin{itemize}
\item \textbf{Stability.} The condition $\eta_t\le 1/(2L)$ guarantees that each local SGD step does not increase the local objective by more than a quadratic term; this yields the standard descent inequality used throughout the proof.
\item \textbf{Effect of Local Steps $E$.} The drift term appears multiplied by $E^2$; thus larger $E$ amplifies the penalty $L E^2\eta_0 D^2/\sqrt{T}$, reflecting the well‑known trade‑off between communication efficiency and convergence speed.
\item \textbf{Comparison to Centralized SGD.} Setting $E=1$ and $K=1$ reduces the bound to the classical $\mathcal O(1/\sqrt{T})$ rate for SGD with variance $\sigma^2/K$.
\item \textbf{Hyper‑parameter Sensitivity.} The bound is linear in the initial stepsize $\eta_0$; choosing $\eta_0\asymp 1/\sqrt{T}$ minimizes the right‑hand side up to constants, reproducing the optimal $\mathcal O(1/\sqrt{T})$ rate.
\end{itemize}

\section*{Preliminaries (Definitions and Lemmas)}
\begin{lemma}[Smoothness Inequality]\label{lem:smooth}
For any $L$‑smooth $f$,
\[
f(w')\le f(w)+\inner{\grad f(w)}{w'-w}+\frac{L}{2}\norm{w'-w}^2 .
\]
\end{lemma}
\begin{lemma}[Unbiased Gradient and Variance]\label{lem:unbiased}
For the stochastic gradient $g_{t}^{i,e}$ defined in \eqref{1},
\[
\E\!\bigl[g_{t}^{i,e}\mid w_{t}^{i,e}\bigr]=\grad f_i\!\bigl(w_{t}^{i,e}\bigr),\qquad
\E\!\bigl[\norm{g_{t}^{i,e}-\grad f_i(w_{t}^{i,e})}^2\mid w_{t}^{i,e}\bigr]\le\sigma^2 .
\]
\end{lemma}
\begin{lemma}[Drift Decomposition]\label{lem:drift}
For any $t$,
\[
\frac1K\sum_{i=1}^K\grad f_i\!\bigl(w_{t}^{i,E}\bigr)
=
\grad f(w_t)+\Delta_t,
\qquad
\Delta_t:=\frac1K\sum_{i=1}^K\bigl[\grad f_i(w_{t}^{i,E})-\grad f_i(w_t)\bigr]+\delta_t .
\]
\end{lemma}
\begin{proof}
Add and subtract $\grad f_i(w_t)$ inside the sum and use the definition of $\delta_t$ in \eqref{3}. \qedhere
\end{proof}

\section*{Main Proof (Step‑by‑Step Derivation)}
We prove Theorem \ref{thm:main}. Throughout the proof we condition on the history up to round $t$ and write $\E_t[\cdot]=\E[\cdot\mid w_0,\dots,w_t]$.

\bigskip
\noindent\textbf{[Step 1] Global Descent after Averaging.}  
From the aggregation rule \eqref{2} and Lemma \ref{lem:smooth} applied to $f$,
\begin{align}
\E_t\!\bigl[f(w_{t+1})\bigr]
&\le f(w_t)
   +\inner{\grad f(w_t)}{\E_t[w_{t+1}-w_t]}
   +\frac{L}{2}\E_t\!\bigl[\norm{w_{t+1}-w_t}^2\bigr].
\tag{5}
\end{align}

\bigskip
\noindent\textbf{[Step 2] Express $w_{t+1}-w_t$ using local updates.}  
Unfolding the $E$ local SGD steps gives
\[
w_{t}^{i,E}=w_t-\eta_t\sum_{e=0}^{E-1} g_{t}^{i,e}.
\]
Thus
\[
w_{t+1}-w_t
= \frac1K\sum_{i=1}^K\bigl(w_{t}^{i,E}-w_t\bigr)
= -\eta_t\frac1K\sum_{i=1}^K\sum_{e=0}^{E-1} g_{t}^{i,e}.
\tag{6}
\]

\bigskip
\noindent\textbf{[Step 3] Conditional expectation of the update.}  
Using Lemma \ref{lem:unbiased},
\[
\E_t\!\bigl[w_{t+1}-w_t\bigr]
= -\eta_t\frac1K\sum_{i=1}^K\sum_{e=0}^{E-1}\grad f_i\!\bigl(w_{t}^{i,e}\bigr).
\tag{7}
\]

\bigskip
\noindent\textbf{[Step 4] Decompose gradients around $w_t$.}  
Add and subtract $\grad f_i(w_t)$ inside the sum:
\begin{align}
\frac1K\sum_{i=1}^K\sum_{e=0}^{E-1}\grad f_i\!\bigl(w_{t}^{i,e}\bigr)
&=
E\grad f(w_t) \;+\; \underbrace{\frac1K\sum_{i=1}^K\sum_{e=0}^{E-1}
\bigl[\grad f_i(w_{t}^{i,e})-\grad f_i(w_t)\bigr]}_{\displaystyle=:\,\Phi_t}
\;+\;E\delta_t .
\tag{8}
\end{align}
Hence
\[
\E_t\!\bigl[w_{t+1}-w_t\bigr]
= -\eta_t\Bigl(E\grad f(w_t)+\Phi_t+E\delta_t\Bigr).
\tag{9}
\]

\bigskip
\noindent\textbf{[Step 5] Plug \eqref{9} into the descent inequality \eqref{5}.}  
Using $\inner{\grad f(w_t)}{-\eta_tE\grad f(w_t)}=-\eta_tE\norm{\grad f(w_t)}^2$ and the Cauchy–Schwarz inequality for the remaining terms,
\begin{align}
\E_t\!\bigl[f(w_{t+1})\bigr]
&\le f(w_t)-\eta_tE\norm{\grad f(w_t)}^2
   +\eta_t\inner{\grad f(w_t)}{-(\Phi_t+E\delta_t)} \nonumber\\
&\qquad
   +\frac{L}{2}\E_t\!\bigl[\norm{w_{t+1}-w_t}^2\bigr].
\tag{10}
\end{align}

\bigskip
\noindent\textbf{[Step 6] Bound the inner product term.}  
By Cauchy–Schwarz and Assumption \ref{as:drift},
\[
\bigl|\inner{\grad f(w_t)}{\Phi_t}\bigr|
\le \norm{\grad f(w_t)}\norm{\Phi_t}, \qquad
\bigl|\inner{\grad f(w_t)}{E\delta_t}\bigr|
\le E\norm{\grad f(w_t)}D .
\tag{11}
\]

\bigskip
\noindent\textbf{[Step 7] Bound $\norm{\Phi_t}$.}  
Each local iterate satisfies, by $L$‑smoothness,
\[
\norm{\grad f_i(w_{t}^{i,e})-\grad f_i(w_t)}\le L\norm{w_{t}^{i,e}-w_t}.
\]
Unrolling the $e$ local updates and using $\eta_t\le 1/(2L)$ yields
\[
\norm{w_{t}^{i,e}-w_t}\le e\eta_t G \le E\eta_t G .
\]
Therefore,
\[
\norm{\Phi_t}
\le \frac1K\sum_{i=1}^K\sum_{e=0}^{E-1} L\,\norm{w_{t}^{i,e}-w_t}
\le L E^2 \eta_t G .
\tag{12}
\]

\bigskip
\noindent\textbf{[Step 8] Bound the second‑moment term $\E_t[\norm{w_{t+1}-w_t}^2]$.}  
From \eqref{6} and independence of stochastic gradients across clients,
\begin{align}
\E_t\!\bigl[\norm{w_{t+1}-w_t}^2\bigr]
&= \eta_t^2 \E_t\!\Bigl[
   \Bigl\|\frac1K\sum_{i=1}^K\sum_{e=0}^{E-1} g_{t}^{i,e}\Bigr\|^2\Bigr] \nonumber\\
&\le \frac{\eta_t^2}{K}\sum_{i=1}^K\sum_{e=0}^{E-1}
   \E_t\!\bigl[\norm{g_{t}^{i,e}}^2\bigr] \qquad\text{(Jensen)} \nonumber\\
&\le \frac{\eta_t^2}{K}\sum_{i=1}^K\sum_{e=0}^{E-1}
   \bigl( \norm{\grad f_i(w_{t}^{i,e})}^2+\sigma^2 \bigr) \nonumber\\
&\le \eta_t^2 E\bigl(G^2+\sigma^2\bigr) .
\tag{13}
\end{align}
Here we used Assumption \ref{as:boundedgrad}.

\bigskip
\noindent\textbf{[Step 9] Assemble the pieces.}  
Insert the bounds \eqref{11}–\eqref{13} into \eqref{10}:
\begin{align}
\E_t\!\bigl[f(w_{t+1})\bigr]
&\le f(w_t)-\eta_tE\norm{\grad f(w_t)}^2
   +\eta_t\norm{\grad f(w_t)}\bigl(L E^2\eta_t G+E D\bigr) \nonumber\\
&\qquad
   +\frac{L}{2}\,\eta_t^2 E\bigl(G^2+\sigma^2\bigr).
\tag{14}
\end{align}
Rearrange terms:
\begin{align}
\eta_tE\norm{\grad f(w_t)}^2
&\le f(w_t)-\E_t\!\bigl[f(w_{t+1})\bigr]
   +\eta_t^2 E\Bigl(L G +\tfrac{L}{2}(G^2+\sigma^2)\Bigr)
   +\eta_t E D \norm{\grad f(w_t)} .
\tag{15}
\end{align}
Apply Young’s inequality $ab\le \frac{1}{2}\bigl(\lambda a^2+\frac{1}{\lambda}b^2\bigr)$ with $\lambda=E\eta_t$ to the mixed term $\eta_t E D \norm{\grad f(w_t)}$:
\[
\eta_t E D \norm{\grad f(w_t)}
\le \frac{\eta_tE}{2}\norm{\grad f(w_t)}^2
   +\frac{\eta_tE D^2}{2}.
\tag{16}
\]
Substituting \eqref{16} into \eqref{15} and cancelling $\frac{\eta_tE}{2}\norm{\grad f(w_t)}^2$ from both sides yields
\begin{align}
\frac{\eta_tE}{2}\norm{\grad f(w_t)}^2
&\le f(w_t)-\E_t\!\bigl[f(w_{t+1})\bigr]
   +\eta_t^2 E C_1
   +\frac{\eta_tE D^2}{2},
\tag{17}
\end{align}
where $C_1:= L G+\frac{L}{2}(G^2+\sigma^2)$ is a constant independent of $t$.

\bigskip
\noindent\textbf{[Step 10] Sum over $t=0,\dots,T-1$.}  
Take total expectation and sum \eqref{17}:
\begin{align}
\sum_{t=0}^{T-1}\frac{\eta_tE}{2}\E\!\bigl[\norm{\grad f(w_t)}^2\bigr]
&\le f(w_0)-\E\!\bigl[f(w_T)\bigr]
   + C_1 E\sum_{t=0}^{T-1}\eta_t^2
   +\frac{E D^2}{2}\sum_{t=0}^{T-1}\eta_t .
\tag{18}
\end{align}
Since $f(w_T)\ge f(w^\star)$,
\[
\sum_{t=0}^{T-1}\eta_t\E\!\bigl[\norm{\grad f(w_t)}^2\bigr]
\le
\frac{2}{E}\bigl(f(w_0)-f(w^\star)\bigr)
   +\frac{2C_1}{E}\sum_{t=0}^{T-1}\eta_t^2
   + D^2\sum_{t=0}^{T-1}\eta_t .
\tag{19}
\]

\bigskip
\noindent\textbf{[Step 11] Relate to suboptimality via convexity.}  
For convex $f$, $\frac{1}{T}\sum_{t=0}^{T-1} f(w_t)-f(w^\star)
\le \frac{1}{T}\sum_{t=0}^{T-1}\inner{\grad f(w_t)}{w_t-w^\star}
\le \frac{1}{T}\sum_{t=0}^{T-1}\norm{\grad f(w_t)}\norm{w_t-w^\star}$.
Using $\norm{w_t-w^\star}\le \norm{w_0-w^\star}+ \sum_{s=0}^{t-1}\norm{w_{s+1}-w_s}$
and the bounded‑gradient assumption, one obtains the standard bound
\[
\E\!\bigl[f(\bar w_T)-f(w^\star)\bigr]
\le \frac{\norm{w_0-w^\star}^2}{2\sum_{t=0}^{T-1}\eta_t}
    +\frac{C_2\sum_{t=0}^{T-1}\eta_t^2}{\sum_{t=0}^{T-1}\eta_t}
    +\frac{D^2\sum_{t=0}^{T-1}\eta_t}{\sum_{t=0}^{T-1}\eta_t},
\tag{20}
\]
where $C_2:=2C_1/E$.
For the stepsize choice $\eta_t=\eta_0/\sqrt{t+1}$ we have
\[
\sum_{t=0}^{T-1}\eta_t \ge 2\eta_0\sqrt{T},\qquad
\sum_{t=0}^{T-1}\eta_t^2 \le \eta_0^2\log T .
\]
Plugging these estimates into \eqref{20} and simplifying gives the explicit bound \eqref{4} stated in Theorem \ref{thm:main}. \qed

\section*{Rate Derivation (With Constants)}
From \eqref{4} we can write
\[
\E\!\bigl[f(\bar w_T)-f(w^\star)\bigr]
\le
\underbrace{\frac{\norm{w_0-w^\star}^2}{2\eta_0}}_{\displaystyle C_{\text{init}}}
\frac{1}{\sqrt{T}}
\;+\;
\underbrace{\frac{L\sigma^2}{K}}_{\displaystyle C_{\sigma}}
\frac{\eta_0}{\sqrt{T}}
\;+\;
\underbrace{L E^2 D^2}_{\displaystyle C_{\delta}}
\frac{\eta_0}{\sqrt{T}} .
\]
Choosing $\eta_0=\Theta(1)$ yields the canonical $\mathcal O(T^{-1/2})$ convergence rate, while the three constants explicitly quantify the influence of initialization, stochastic variance, and client drift.

\section*{Special Cases}
\begin{enumerate}[label=\alph*.]
\item \textbf{No client drift ($D=0$).} The bound reduces to the standard SGD rate with variance reduced by the factor $K$ (because the average of $K$ unbiased gradients has variance $\sigma^2/K$).
\item \textbf{Single local step ($E=1$).} The drift penalty becomes $L\eta_0 D^2/\sqrt{T}$, matching the analysis of \emph{parallel} SGD.
\item \textbf{Homogeneous data ($\delta_t\equiv0$).} The term $D$ vanishes; the algorithm behaves like centralized SGD with an effective stepsize multiplied by $E$.
\end{enumerate}

\end{document}